% Lösungen Einheit 4 von 5 – Dreiecke konstruieren (zusammenbau v0.1)
\einheitenkopf{Einheit 4 von 5 · Dreiecke konstruieren}

\erg{20}{a) SWS \quad b) c = 7 cm von A nach B; Kreisbogen um A mit 5 cm, Kreisbogen um B mit 6 cm; der Schnittpunkt ist C. Kontrolle: $\gamma \approx 78^\circ$ \quad c) c = 6 cm von A nach B; in A den Winkel $50^\circ$ antragen; auf dem freien Schenkel b = 5 cm abtragen, Endpunkt C. Kontrolle: a ≈ 4,7 cm \quad d) c = 7 cm von A nach B; in A $38^\circ$, in B $71^\circ$ antragen; die freien Schenkel schneiden sich in C. Kontrolle: $\gamma = 71^\circ$ \quad e) c = 5 cm von A nach B; in B den Winkel $64^\circ$ antragen; Kreisbogen um A mit 6 cm schneidet den freien Schenkel in C. Kontrolle: a ≈ 6,2 cm \quad f) c = 6 cm von A nach B; Kreisbögen um A und um B mit je 5 cm; der Schnittpunkt ist C. Kontrolle: $\gamma \approx 74^\circ$ \quad g) in C einen rechten Winkel antragen; auf den Schenkeln a = 4 cm (Endpunkt B) und b = 6 cm (Endpunkt A) abtragen; A mit B verbinden. Kontrolle: c ≈ 7,2 cm \quad h) Erst: Seite c zeichnen, Endpunkte A und B. Dann: In A den Winkel α an c antragen. Dann: Auf dem freien Schenkel von A aus b abtragen, Endpunkt C. Zuletzt: C mit B verbinden. \quad i) nein: $4 + 5 = 9$ ist kleiner als 10 \quad j) A ↔ R, B ↔ P, C ↔ Q \quad k) $b + c > 12$ m}

20b)

\dreieck{(0,0)}{(7,0)}{(2.71,4.20)}{a}{b}{c}{\alpha}{\beta}{\gamma}


20c)

\dreieck{(0,0)}{(6,0)}{(3.21,3.83)}{a}{b}{c}{\alpha}{\beta}{\gamma}


20d)

\dreieck{(0,0)}{(7,0)}{(5.52,4.31)}{a}{b}{c}{\alpha}{\beta}{\gamma}


20e)

\dreieck{(0,0)}{(5,0)}{(2.30,5.54)}{a}{b}{c}{\alpha}{\beta}{\gamma}


20f)

\dreieck{(0,0)}{(6,0)}{(3.00,4.00)}{a}{b}{c}{\alpha}{\beta}{\gamma}


20g)

\dreieckrw{4}{6}{a}{b}{c}

\erg{21}{a) Ecken A, B, C gegen den Uhrzeigersinn, a gegenüber von A, b gegenüber von B, c gegenüber von C; markiert: b, c und der Winkel α bei A zwischen ihnen}

21a)

\dreieck{(0,0)}{(5,0)}{(1.5,3)}{a}{b}{c}{\alpha}{\beta}{\gamma}

\erg{22}{a) Fehler: WWW ist kein Kongruenzsatz; drei Winkel legen nur die Form fest, nicht die Größe – es gibt beliebig viele verschieden große Dreiecke mit diesen Winkeln. \quad b) Ein Dreieck lässt sich vergrößern oder verkleinern, ohne dass sich seine Winkel ändern; drei Winkel legen die Größe nicht fest. \quad c) Nein: 2,0 m + 1,2 m = 3,2 m ist kürzer als 3,5 m, die Bretter schließen kein Dreieck.}
